\bookchapter{الفصل 66: المشروع الثاني — تطبيق الاستعلام عن الطقس (نسخة شبكية)}

\chapterintro{دفتر الحسابات كان «تدريبًا على القوة الداخلية»، وتطبيق الطقس «تعرفًا على العالم» — فهو يتعامل مع الإنترنت. في هذا المشروع ستتعلم: طلب API مجاني، وتحليل JSON، والتعامل مع حالات الواجهة الثلاث «جار التحميل / النجاح / الفشل»، وكيف تخفي مفتاح API جيدًا. وبإتمامه ستملك القدرة على صنع أي «تطبيق شبكي» — الأخبار وأسعار الصرف والترجمة، فالطريقة واحدة عينها.}

\booksection{أولًا: طلب واجهة طقس مجانية (API)}

للتدريب الشخصي ننصح بـ OpenWeatherMap أو الخدمة الصينية QWeather (طقس هفنغ) — فكلاهما يمنح حصة مجانية تعمل فور التسجيل. ولنأخذ OpenWeatherMap مثالًا:

\begin{enumerate}
  \item سجل حسابًا في صفحة API keys أنشئ مفتاحًا (apikey)؛
  \item النسخة المجانية تستدعي واجهة Current Weather: api.openweathermap.org/data/2.5/weather? q=beijing\&appid=مفتاحك\&units=metric؛
  \item جربها أولًا في المتصفح وانظر كيف يبدو JSON العائد — فهذا «الواجب التمهيدي» الذي لا بد منه قبل ربط أي API.
\end{enumerate}

\textbf{تنبيه لتجنب الفخاخ: المفتاح لا يجوز أن يكتب ثابتًا في الكود أبدًا}\par

مفتاح API أشبه بـ«كلمة سر محفظتك»؛ فما إن يكتب في الكود ويرفع إلى GitHub حتى تجرفه الزواحف خلال ساعات قليلة ويستخدم بغير حق. والطريقة الصحيحة في iOS: ضعه في Config.xcconfig (ملف إعدادات يتجاهله gitignore)، ثم اقرأه في الكود عبر Info.plist. ابحث عن «إدارة مفاتيح iOS باستخدام xcconfig» وعدّل وفق ما تجد.

\booksection{ثانيًا: طبقة الشبكة — URLSession مع Codable}

عرّف النموذج أولًا وفق JSON الذي يعيده الـ API، ثم اكتب الخدمة:

\begin{lstlisting}[language=Swift]
// JSON: {"weather":[{"description":"صافٍ"}],"main":{"temp":26.5,"humidity":40},"name":"الرياض"}
struct WeatherResponse: Codable {
    let weather: [Weather]
    let main: Main
    let name: String
}
struct Weather: Codable { let description: String }
struct Main: Codable { let temp: Double; let humidity: Int }

struct WeatherService {
    let key: String   // يُقرأ من xcconfig

    func fetchWeather(city: String) async throws -> WeatherResponse {
        var comps = URLComponents(string: "https://api.openweathermap.org/data/2.5/weather")!
        comps.queryItems = [
            URLQueryItem(name: "q", value: city),
            URLQueryItem(name: "appid", value: key),
            URLQueryItem(name: "units", value: "metric"),
            URLQueryItem(name: "lang", value: "zh_cn")
        ]
        let (data, _) = try await URLSession.shared.data(from: comps.url!)
        return try JSONDecoder().decode(WeatherResponse.self, from: data)
    }
}
\end{lstlisting}

\booksection{ثالثًا: حالات الواجهة الثلاث — جار التحميل / النجاح / الفشل}

روح التطبيق الشبكي في «إدارة الحالة». بطء الشبكة، وانقطاعها، وخطأ في اسم المدينة… لكل حالة يجب أن توضح الواجهة للمستخدم ما يجري. عبّر عنها بأناقة بـ enum مع switch:

\begin{lstlisting}[language=Swift]
enum LoadState {
    case loading                  // جارٍ التحميل
    case success(WeatherResponse) // نجاح
    case error(String)            // فشل
}

@Observable
class WeatherViewModel {
    var state: LoadState = .loading
    private let service = WeatherService(key: Config.apiKey)

    func search(city: String) async {
        state = .loading
        do {
            let data = try await service.fetchWeather(city: city)
            state = .success(data)
        } catch {
            state = .error("تعذر العثور على هذه المدينة، جرّب اسمًا آخر؟")
        }
    }
}

// في الواجهة:
switch vm.state {
case .loading: ProgressView("جارٍ التحميل…")
case .success(let d): WeatherContent(data: d)
case .error(let msg): Text(msg).foregroundColor(.secondary)
}
\end{lstlisting}

هذا النمط الثلاثي loading/success/error هو «القالب المضمون متعدد الأغراض» للواجهات الشبكية؛ احفظه الآن، وسترسم على منواله كل تطبيقاتك الشبكية لاحقًا.

\booksection{رابعًا: كود الواجهة عمليًا — نرسم الطقس}

\begin{lstlisting}[language=Swift]
struct WeatherContent: View {
    let data: WeatherResponse

    var body: some View {
        VStack(spacing: 12) {
            Text(data.name)
                .font(.title.bold())
            Text("\(Int(data.main.temp))°C")
                .font(.system(size: 76, weight: .thin))
            Text(data.weather.first?.description ?? "غير معروف")
                .font(.title3)
                .foregroundColor(.secondary)
            HStack(spacing: 24) {
                Label("الرطوبة \(data.main.humidity)%", systemImage: "humidity")
                Label("الإحساس \(Int(data.main.temp))°C", systemImage: "thermometer")
            }
        }
        .padding()
    }
}
\end{lstlisting}

تفاصيل صغيرة تستحق التعلم: first? لالتقاط العنصر الأول بأمان (قد تكون القائمة فارغة)، و?? “غير معروف” نصًا احتياطيًا؛ وعرض درجة الحرارة بخط رقيق حجمه 76 وصفة كلاسيكية في تطبيقات الطقس؛ وLabel هو المكون الرسمي «أيقونة + نص»، وهو أنظم من تركيب HStack يدويًا.

\booksection{خامسًا: تفاصيل التفاعل والمدن المفضلة}

\begin{itemize}
  \item حقل البحث: المُعدِّل .searchable في الأعلى يجلب تفاعل البحث من النظام، وعند الإرسال يستدعي vm.search؛
  \item المدن المفضلة: خزن المدن التي تتابعها كثيرًا في UserDefaults، فيستعلم عنها تلقائيًا عند الفتح التالي؛
  \item السحب للتحديث: .refreshable \{ await vm.search(...) \} تكفي سطرًا واحدًا؛
  \item أيقونة الطقس: استخدم Image(systemName:) بالربط حسب الوصف (sun.max / cloud.rain / snowflake)، فتنبض الواجهة بالحياة فورًا.
\end{itemize}

\begin{tcbinfo}{نصيحة: «عويصة» أسماء المدن الصينية}
دعم OpenWeatherMap لأسماء المدن بالحروف الصينية متوسط، فالبحث عن «بكين» بالحروف الصينية قد يعود بلا نتائج بينما “beijing” يعمل عاديًا. والحل الاحترافي هو تضمين جدول تحويل «اسم المدينة اسم إنجليزي/إحداثيات»، أو الربط بخدمة طقس محلية. في مرحلة التدريب يكفي أن تدعم أسماء المدن بالإنجليزية، فلا تشن حربًا مع هذا القصور.

\end{tcbinfo}

\begin{tcbpractice}{هيا نجرّب: قائمة قبول تطبيق الطقس}
\booksection{سادسًا: كيف تنقح طلبات الشبكة — الطباعة والتقاط الحزم ومحاكاة الشبكة البطيئة}

إذا اختل تطبيق شبكي ففي تسع حالات من عشر يكون السبب «طلبًا لم يرسل صحيحًا» أو «استجابة لم تحلل صحيحة». ثلاث وسائل للتنقيح:

\begin{enumerate}
  \item طباعة البيانات الخام: اطبع String(data: data, encoding: .utf8) لترى JSON الخام — فتتبين أسماء الحقول وأنواعها من نظرة واحدة، وهو أسرع من التخمين بعشرة أضعاف؛
  \item أدوات التقاط الحزم: Proxyman المجاني على Mac أو Charles مفتوح المصدر يعترض طلبات التطبيق واستجاباته كلها، فترى طبقة الشبكة بوضوح أشعة X؛
  \item محاكاة الشبكة الضعيفة: من قائمة المحاكي Features Network Link Conditioner (تتطلب تنزيل أداة من موقع المطورين)، تحاكي شبكة 3G والتأخر العالي — فتكشف مشكلات تجربة «الدوامة التي لا تنتهي».
\end{enumerate}

\begin{tcbinfo}{نصيحة: هل رفض ATS طلبك؟}
إذا ظهر خطأ “App Transport Security” بمجرد إرسال الطلب فالأرجح أنك استخدمت عنوان http:// نصيًا صريحًا. الواجهات الرسمية يجب أن تكون HTTPS؛ وفي التنقيح المحلي افتح استثناء مؤقت (أضف NSAllowsLocalNetworking إلى Info.plist)، وتذكر تنظيفه قبل النشر.

\end{tcbinfo}

(1) البحث عن الطقس يعمل بشكل طبيعي؛ (2) عند قطع الشبكة يعرض التطبيق رسالة خطأ أنيقة دون انهيار (جرب بإطفاء شبكة Mac من داخل المحاكي!)؛ (3) إدخال مدينة غير موجودة يعطي تنبيهًا ودودًا؛ (4) حفظ مدينتين مفضلتين يبقيهما بعد إعادة تشغيل التطبيق؛ (5) الواجهة لا تؤذي العين في الوضع الداكن. وإذا اجتازت البنود الخمسة كلها يكون هذا المشروع قد تخرج.

\end{tcbpractice}

\begin{tcbsummary}{خلاصة الفصل}
\begin{itemize}
  \item جرب الـ API في المتصفح وافهم JSON قبل كتابة الكود — فـ«الاستطلاع المسبق» لا يخسر أبدًا
  \item أخف مفتاح API في xcconfig، ولا يدخل مستودع الكود أبدًا
  \item نموذج Codable مع URLSession وasync/await: إذا تطابقت الحقول تحلل JSON تلقائيًا
  \item نمط الحالات الثلاث loading/success/error هو قالب كل الأغراض للواجهات الشبكية
  \item انقطاع الشبكة، والإدخال الفارغ، ومدينة غير موجودة — معالجة الأخطاء هي حد النجاح للتطبيقات الشبكية
\end{itemize}

\end{tcbsummary}
